\section{Point-cloud track: PET on the raw reconstructed clusters}
\label{sec:pet}

The full promise of OmniFold is unfolding on the event's \emph{raw}
representation rather than on engineered scalar summaries. This section
documents that track: a point-cloud transformer (PET, the architecture of
OmniLearn~\cite{Mikuni:2024qsr}, vendored with the upstream OmniFold
implementation) unfolds the same ME~FHC data using, on the reconstructed
side, the actual calorimetric recoil clusters --- per-cluster energy,
transverse position and $z$ (\texttt{cluster\_energy}, \texttt{cluster\_pos},
\texttt{cluster\_z}, filtered on \texttt{cluster\_isMuontrack} via
\texttt{CVUniverse::GetRecoClusters()}) alongside the muon kinematics. An
earlier attempt that read the wrong, mostly-empty
\texttt{ExtraEnergyClusters\_*} branches was discarded and the chain rebuilt
from the corrected dump. \textbf{Maturity:} this is a validated \emph{method
milestone} --- a real, absolutely-normalised, full-statistics cross section
with its own closure gate and a first systematic budget --- but it is not yet
at the production GBDT result's maturity (its training is
subsample-limited; see below).

\subsection{From shape check to absolute cross section}
\label{sec:pet-abs}
The PET network trains on a GPU on a 2M-event subsample (2 OmniFold
iterations); the learned step-2 truth reweight is then applied to the
\emph{full} 32.8M-event generator sample
(\texttt{MultiFold.reweight}), binned on the 4D grid, and pushed through the
standard extraction (\texttt{extract\_cross\_section\_nd} with the GBDT
run's completeness --- completeness is reweight-independent --- and the same
flux/POT/nucleon constants). Three gates validate the absolute path:
\begin{enumerate}[itemsep=2pt]
  \item \textbf{Full-statistics reweight sanity}: the pushed weights over the
        32.8M generator events are finite with mean $1.028$;
  \item \textbf{Closure (decisive)}: unfolding MC-reco-as-pseudo-data
        recovers MC truth to total $0.9884$ with uniform per-axis median
        residuals of $\sim\SI{1.1}{\percent}$ --- the absolute machinery is
        unbiased;
  \item \textbf{Cross-check vs GBDT}: on data, the PET absolute total is
        $\SI{2.796e-38}{cm^2/nucleon}$ versus the GBDT
        $\SI{3.066e-38}{}$, ratio $\petRatio$, with per-axis median differences
        of \SIrange{6.5}{9.9}{\percent} (Fig.~\ref{fig:petabs}).
\end{enumerate}
Because the closure is exact to $\sim\SI{\petClosure}{\percent}$, the $\sim\SI{9}{\percent}$
data-side gap is a genuine \emph{training-configuration} difference (PET
trained on the 2M subsample with 2 iterations versus the full-statistics
5-iteration GBDT), not a normalisation bug: an under-iterated unfold pushes
the result less far from the prior. A higher-epoch retrain moved the total by
$-1.6\%$ (ratio $0.897$), confirming the training-config sensitivity; closing
the gap toward the ML band needs more events/iterations and is deferred.

\begin{figure}[H]
  \centering
  \includegraphics[width=0.95\textwidth]{pet_vs_gbdt_absolute.png}
  \caption{Absolutely-normalised 4D cross section from the PET point-cloud
  unfold (full-statistics reweight of the generator sample) versus the
  production scalar GBDT result, projected on each axis. The PET/GBDT total
  ratio is $\petRatio$; the closure gate (text) shows the absolute machinery
  itself is unbiased to $\sim\SI{\petClosure}{\percent}$.}
  \label{fig:petabs}
\end{figure}

\subsection{PET systematic budget (frozen reweighter)}
\label{sec:pet-syst}
A first full uncertainty budget is attached with the \emph{frozen-reweighter}
method: the PET step-2 weights are held fixed while each systematic
universe's truth reweight ratios re-bin the generator sample --- 13 knob
bands, the 100-universe flux multisim, a 100-replica Poisson bootstrap
(statistical), an alternative-training difference (ML), and the detector
(lateral) bands transferred from the GBDT 4D covariance by fractional
projection. The combined budget on the PET reported bins has median
per-bin fractions of \SI{8.2}{\percent} (systematic), \SI{4.2}{\percent}
(statistical), \SI{3.3}{\percent} (ML), \SI{4.0}{\percent} (lateral,
transferred), totalling $\approx\SI{12.3}{\percent}$ --- comparable to the
GBDT 4D budget's \SI{13.5}{\percent}. The reweight bank underlying the
systematic block is built from the truth-only miss-row dump, with per-event
alignment verified bit-identical over all \num{32.8}~million rows and the
bank-independent statistical and ML blocks reproducing as a control. The
residual gap to a fully self-contained PET budget is
per-lateral re-inference on shifted reconstructed clouds, which requires
re-running the PET network per lateral universe and is deferred.

\subsection{What the PET track demonstrates}
The point-cloud track closes the methodological loop begun by the scalar
GBDT/MLP cross-checks (\S\ref{sec:validation}): the unfolded physics is
stable across classifier family (GBDT, MLP, transformer) \emph{and} input
representation (engineered scalars, raw reconstructed clusters), now
including the absolute normalisation. It also establishes the
infrastructure --- point-cloud DataLoader, GPU training, full-sample
reweighting, frozen-reweighter budgets --- for unfolding directly on
low-level MINERvA data, the natural endpoint of the unbinned program.
